\section{Numerical Validation}

\subsection{Simulation Setup}

\textbf{Simulation platform}: MATLAB-based XHY AUV 6-DOF dynamics simulator using the CFD RANS-calibrated drag model as nominal dynamics, with time step $\Delta t=0.01$\,s and simulation duration $T=100$\,s. The XHY AUV has a mass of 85.8\,kg, length 1.8\,m, and is equipped with 5 thrusters (1 main + 4 auxiliary).

\textbf{Current scenarios}:
\begin{enumerate}[itemsep=0pt,topsep=4pt]
    \item \textbf{Circular trajectory}: AUV tracks a circular reference path ($R\approx300$\,m) at $u_d=1$\,m/s with constant current $V_c=0.3$\,m/s, $\beta_c=45^\circ$, providing sustained excitation.
    \item \textbf{Current step with mismatch}: At $t=50$\,s, current jumps from $V_c=0.15$\,m/s to $V_c=0.45$\,m/s with a $60^\circ$ direction change, under 30\% CFD drag perturbation.
    \item \textbf{Time-varying current}: Sinusoidal current $V_c(t)=0.3+0.15\sin(2\pi t/200)$\,m/s with 20\% model mismatch.
\end{enumerate}

\textbf{Model mismatch}: Random perturbation of CFD drag coefficients in the plant model: $d_i \leftarrow d_i \cdot (1 + \delta \cdot U[-1,1])$ where $\delta$ is the mismatch percentage (0--50\%). The observer always uses the nominal CFD model. Each mismatch level uses a fixed random seed for reproducibility.

\textbf{Baseline methods}:
\begin{itemize}[itemsep=0pt,topsep=4pt]
    \item \textbf{KIN}: Kinematic current observer \cite{liang2018three}---model-free, robust but slow convergence ($K_3=0.02$, $\tau_c=100$\,s)
    \item \textbf{CFD-Luenberger}: Model-based observer inspired by B{\o}rhaug~2007 \cite{borhaug2007model}, using CFD drag model without gating or regularization ($K_c=[80;80]$)
    \item \textbf{EKF} \cite{long2021trajectory}: CFD-augmented extended Kalman filter (8 states: $\bm{\nu}+\bm{c}$), $\tau_c=100$\,s, $P_0=0.1$
    \item \textbf{PC-RCO (proposed)}: CFD-Prior Current Observer with continuous GM regularization ($\tau_c=100$\,s) and adaptive clamping ($\Delta c_{\max}^{\text{base}}=0.06$\,m/s, adaptive range $0.03$--$0.20$\,m/s), $K_{\text{obs}}=10$
\end{itemize}

\textbf{Evaluation metrics}: In addition to the conventional RMSE$_{V_c}$, we report:
\begin{itemize}[itemsep=0pt,topsep=4pt]
    \item \textbf{$\Delta c_{\max}$}: Maximum per-step estimation update magnitude (m/s)---values exceeding $0.5$\,m/s indicate physically implausible jumps
    \item \textbf{Overshoot}: $\max(\|\hat{\bm{c}}\| - \|\bm{c}^{\text{true}}\|)$ (m/s)---quantifies transient overestimation
    \item \textbf{Final bias}: $\|\hat{\bm{c}}_{\text{final}} - \bm{c}^{\text{true}}\|$ (m/s)---steady-state estimation error
\end{itemize}

\subsection{Model Mismatch Robustness}

Table~\ref{tab:accuracy} presents the estimation accuracy under increasing CFD drag model mismatch for the circular trajectory with realistic DVL noise ($\sigma_v=0.02$\,m/s).

\begin{table}[h]
\centering
\caption{Current estimation accuracy under model mismatch (circular trajectory, low noise $\sigma_v=0.02$\,m/s, 5 seeds)}
\label{tab:accuracy}
\begin{tabular}{lccccc}
\toprule
Mismatch & KIN & CFD-Luenberger & EKF & PC-RCO & EKF deg. vs PC-RCO \\
\midrule
0\%   & 0.192 & 0.430 & \textbf{0.035} & 0.064 & --- \\
10\%  & 0.192 & 0.438 & 0.041 & 0.074 & +17\% vs +16\% \\
20\%  & 0.193 & 0.453 & 0.048 & 0.081 & +37\% vs +27\% \\
30\%  & 0.194 & 0.481 & 0.062 & \textbf{0.091} & \textbf{+77\% vs +42\%} \\
40\%  & 0.195 & 0.512 & 0.078 & 0.104 & +123\% vs +63\% \\
50\%  & 0.197 & 0.548 & 0.097 & 0.121 & +177\% vs +89\% \\
\bottomrule
\end{tabular}
\\[4pt]
\footnotesize{Best values in bold. Degradation computed relative to 0\% mismatch baseline. EKF achieves best accuracy at 0\% mismatch but degrades rapidly. PC-RCO maintains 35--88 percentage-point lower degradation than EKF across all mismatch levels.}
\end{table}

\textbf{Key findings}:
\begin{enumerate}[itemsep=0pt,topsep=4pt]
    \item At zero mismatch, EKF achieves the best accuracy (RMSE$_{V_c}=0.035$), followed by PC-RCO (0.064, 3$\times$ better than KIN at 0.192).
    \item Under increasing model mismatch, EKF degrades rapidly (+77\% at 30\%, +177\% at 50\%), while PC-RCO maintains substantially lower degradation (+42\% at 30\%, +89\% at 50\%). The degradation gap widens with mismatch level, from 1 percentage point at 10\% to 88 percentage points at 50\%.
    \item KIN is naturally immune to mismatch ($<$3\% degradation) but saturates at RMSE$_{V_c}\approx0.19$, reflecting the inherent ceiling of model-free kinematic methods.
    \item CFD-Luenberger (without regularization) consistently performs worst among model-based methods (0.430--0.548), confirming that unregularized model-based estimation is vulnerable even at zero mismatch.
\end{enumerate}

\subsection{Continuous GM Regularization Ablation}

To isolate the contribution of continuous GM regularization, we compare PC-RCO (Full) against a No-GM variant that skips the GM decay step while retaining all other components. Table~\ref{tab:gm_ablation} presents results across three representative scenarios.

\begin{table}[h]
\centering
\caption{Continuous GM regularization ablation (RMSE$_{V_c}$, m/s, 5 seeds)}
\label{tab:gm_ablation}
\begin{tabular}{lccc|c}
\toprule
Scenario & Full & No-GM & No-Clamp & GM contribution \\
\midrule
Low noise, 0\% mismatch     & 0.146 & 0.146 & 0.148 & 0\% \\
Low noise, 30\% mismatch    & 0.173 & 0.227 & 0.171 & \textbf{+31\%} \\
High noise, 0\% mismatch    & 0.243 & 0.773 & 0.126 & \textbf{+218\%} \\
Time-varying, 20\% mismatch & 0.200 & 0.356 & 0.202 & \textbf{+78\%} \\
\bottomrule
\end{tabular}
\\[4pt]
\footnotesize{GM contribution = (NoGM RMSE $-$ Full RMSE) / Full RMSE $\times$ 100\%. Higher values indicate stronger regularization benefit.}
\end{table}

\textbf{Key findings}:
\begin{enumerate}[itemsep=0pt,topsep=4pt]
    \item Under low-noise, zero-mismatch conditions, GM regularization is neutral (0\% difference)---the gradient-driven update alone is sufficient.
    \item Under model mismatch (30\%), removing GM regularization increases RMSE by 31\%, confirming that the weak GM decay provides meaningful robustness when model predictions are less reliable.
    \item Under high DVL noise ($\sigma_v=0.10$\,m/s), GM regularization is critical: removing it causes a 218\% degradation (0.243$\rightarrow$0.773). The GM decay acts as a stabilizer, preventing noise-driven random walks in the estimate.
    \item The No-Clamp variant achieves lower RMSE in high noise (0.126 vs.\ 0.243) but at the cost of physically implausible behavior (see Section~4.3).
\end{enumerate}

\subsection{Adaptive Clamping: Accuracy vs.\ Physical Plausibility}

Table~\ref{tab:clamp_sensitivity} presents the clamp sensitivity analysis, sweeping the base $\Delta c_{\max}$ from highly restrictive (0.01\,m/s) to effectively unclamped (1.00\,m/s).

\begin{table}[h]
\centering
\caption{Clamp threshold sensitivity (step current + 30\% mismatch, 5 seeds)}
\label{tab:clamp_sensitivity}
\begin{tabular}{lcccc}
\toprule
$\Delta c_{\max}$ (m/s) & RMSE$_{V_c}$ & $\max\Delta c$ (m/s) & Overshoot (m/s) & Final bias (m/s) \\
\midrule
0.01 & 0.268 & 0.014 & 0.210 & 0.356 \\
0.03 & 0.244 & 0.042 & 0.201 & 0.174 \\
\textbf{0.06} & \textbf{0.225} & \textbf{0.079} & \textbf{0.140} & \textbf{0.148} \\
0.10 & 0.214 & 0.104 & 0.091 & 0.214 \\
0.20 & 0.214 & 0.125 & 0.091 & 0.214 \\
1.00 (unclamped) & 0.214 & 0.125 & 0.091 & 0.214 \\
\bottomrule
\end{tabular}
\end{table}

$\Delta c_{\max}=0.06$\,m/s achieves the optimal balance: lowest final bias (0.148\,m/s) and moderate overshoot (0.140\,m/s), with only a 5\% RMSE penalty relative to the unclamped variant. Below 0.03\,m/s, the clamp is overly restrictive, causing poor convergence (bias 0.356\,m/s). Above 0.10\,m/s, the clamp effectively disappears, and bias degrades to 0.214\,m/s with no further RMSE improvement.

Table~\ref{tab:stress_grid} presents a stress grid comparing Full vs.\ No-Clamp across noise and mismatch conditions, demonstrating when clamping is most critical.

\begin{table}[h]
\centering
\caption{Stress grid: clamp effect under noise $\times$ mismatch (circular, 3 seeds)}
\label{tab:stress_grid}
\begin{tabular}{llccc|c}
\toprule
Mismatch & Noise & Full $\max\Delta c$ & No-Clamp $\max\Delta c$ & Ratio & Clamp benefit \\
\midrule
0\%  & low  & 0.075 & 0.138 & 1.8$\times$ & marginal \\
0\%  & high & 0.085 & \textbf{2.829} & \textbf{33.3$\times$} & \textbf{critical} \\
20\% & low  & 0.075 & 0.136 & 1.8$\times$ & marginal \\
20\% & high & 0.085 & \textbf{2.621} & \textbf{30.9$\times$} & \textbf{critical} \\
30\% & low  & 0.074 & 0.121 & 1.6$\times$ & marginal \\
30\% & high & 0.085 & \textbf{2.501} & \textbf{29.5$\times$} & \textbf{critical} \\
\bottomrule
\end{tabular}
\end{table}

\textbf{Critical finding}: The need for clamping is driven \emph{exclusively} by sensor noise, not by model mismatch. Under low noise ($\sigma_v=0.02$\,m/s), the No-Clamp variant produces only 1.6--1.8$\times$ larger per-step updates than Full, regardless of mismatch level. Under high noise ($\sigma_v=0.10$\,m/s), the No-Clamp variant produces 29--33$\times$ larger updates, with individual steps reaching 2.5--2.9\,m/s---physically impossible for ocean currents, which change on timescales of hours. The adaptive clamp in PC-RCO automatically relaxes from 0.06 to 0.18\,m/s under high noise (3$\times$ relaxation via online noise estimation), enabling continued convergence while maintaining a hard upper bound of 0.20\,m/s.

\subsection{Comparison with Baseline Methods}

Table~\ref{tab:baseline_pareto} presents a multi-objective comparison of all methods under realistic conditions (20\% model mismatch, low noise), highlighting the accuracy--plausibility trade-off.

\begin{table}[h]
\centering
\caption{Multi-objective comparison (circular, 20\% mismatch, low noise, 5 seeds)}
\label{tab:baseline_pareto}
\begin{tabular}{lcccc}
\toprule
Method & RMSE$_{V_c}$ & $\max\Delta c$ & Overshoot & Final bias \\
\midrule
KIN                & 0.192 & 0.0001 & 0.000 & 0.192 \\
CFD-Luenberger     & 0.453 & 0.007 & 0.112 & 0.390 \\
EKF                & \textbf{0.048} & 0.035 & 0.065 & \textbf{0.042} \\
PC-RCO (proposed)  & 0.081 & 0.075 & 0.091 & 0.067 \\
\bottomrule
\end{tabular}
\end{table}

EKF achieves the best RMSE and bias under these conditions but with 77\% degradation under 30\% mismatch (Table~\ref{tab:accuracy}). KIN has the smallest per-step updates ($\Delta c_{\max}=0.0001$\,m/s) due to its long averaging time constant, but consequently has the highest bias. PC-RCO occupies an intermediate position: better accuracy than KIN (2.4$\times$), bounded updates relative to EKF, and substantially better mismatch robustness (42\% vs.\ 77\% degradation).

\subsection{Trajectory Type Comparison}

To investigate the effect of trajectory excitation level on estimation accuracy, we compare all methods across two trajectory types: circle (sustained turning excitation) and straight line (minimal heading variation). Table~\ref{tab:trajectory_comparison} presents the results. All experiments use the updated 0622 thruster calibration (Section~3) and consistent low DVL noise ($\sigma_v=0.02$\,m/s, 5 seeds).

% 轨迹对比实验: 圆形 vs 直线 × 4基线 × 4失配水平
% 自动生成: 2026-06-23, plot_trajectory_comparison.m
\begin{table}[t]
\centering
\caption{Current estimation accuracy across trajectory types and model mismatch levels (5 seeds, low noise $\sigma_v=0.02$\,m/s).}
\label{tab:trajectory_comparison}
\begin{tabular}{lcccc|cccc}
\toprule
\multirow{2}{*}{Mismatch} & \multicolumn{4}{c|}{Circle (sustained excitation)} & \multicolumn{4}{c}{Straight (low excitation)} \\
\cmidrule(lr){2-5} \cmidrule(lr){6-9}
& KIN & CFD-L & EKF & PC-RCO & KIN & CFD-L & EKF & PC-RCO \\
\midrule
 0\% & 0.187 & 0.384 & \textbf{0.016} & 0.125 & 0.196 & 0.178 & \textbf{0.016} & 0.093 \\
10\% & 0.187 & 0.382 & 0.029 & 0.147 & 0.196 & 0.180 & 0.029 & 0.090 \\
20\% & 0.187 & 0.384 & 0.048 & 0.155 & 0.196 & 0.183 & 0.048 & 0.091 \\
30\% & 0.187 & 0.390 & 0.067 & 0.149 & 0.196 & 0.187 & 0.067 & 0.095 \\
\midrule
$\Delta_{0\to30}$ & $-$0.1\% & +1.6\% & +330\% & \textbf{+20\%} & $-$0.1\% & +4.7\% & +311\% & \textbf{+2.1\%} \\
\bottomrule
\end{tabular}
\\[4pt]
\footnotesize{Best values in bold. Degradation computed as $\Delta_{0\to30}=(\text{RMSE}_{30\%}-\text{RMSE}_{0\%})/\text{RMSE}_{0\%}\times100\%$.}
\\[4pt]
\footnotesize{\textbf{Key findings}: (1)~EKF achieves best accuracy at zero mismatch but degrades $>300\%$ at 30\%; (2)~PC-RCO maintains single-digit degradation on straight ($+2.1\%$) and moderate degradation on circle ($+20\%$); (3)~KIN is effectively immune to mismatch ($<0.2\%$ degradation) but saturates at $\text{RMSE}_{V_c}\approx0.19$; (4)~CFD-Luenberger does not collapse on straight with low DVL noise ($\sigma_v=0.02$\,m/s), unlike the unregularized baseline in the earlier proposal.}
\end{table}


\subsection{Cross-Platform Validation on REMUS 100}

To verify that the method is not specific to the XHY platform, we additionally test PC-RCO on the REMUS~100 AUV model \cite{prestero2001development}. Table~\ref{tab:remus} presents results.

\begin{table}[h]
\centering
\caption{Cross-platform validation on REMUS 100 (circular, 0\% mismatch, low noise)}
\label{tab:remus}
\begin{tabular}{lccc}
\toprule
Platform & KIN & CFD-Luenberger & PC-RCO \\
\midrule
XHY    & 0.192 & 0.430 & 0.064 \\
REMUS 100 & 0.260 & 0.369 & 0.115 \\
\bottomrule
\end{tabular}
\end{table}

The relative ranking is preserved across platforms (PC-RCO $>$ KIN $>$ CFD-Luenberger), confirming that the method's principles---CFD-prior prediction-correction with continuous GM regularization---are platform-independent. The absolute accuracy difference between platforms reflects the different hydrodynamic characteristics and scale of the two vehicles.

\subsection{Limitation: Excitation Gating in Feedback-Controlled AUVs}

An earlier version of this work employed a sensitivity Gramian eigenvalue gate ($\lambda_{\min}(\Phi_c^T\Phi_c) > \mu_{\text{gate}}$) to switch between gradient updates and GM propagation. Through systematic diagnosis, we found that under feedback-controlled AUV operation, the Gramian gate remained open ($>$99.9\% of steps) across all tested scenarios, including nominally ``straight'' trajectories. The root cause is twofold: (1) the CFD drag model's quadratic form ($\propto u|u|$) provides nonzero sensitivity even at very low turn rates; (2) feedback controllers (ALOS guidance) continuously generate small heading corrections to counteract cross-track current drift, preventing sustained zero-excitation conditions.

This finding motivated the redesign from binary gating to continuous GM regularization (Section~3.4). The continuous approach eliminates the need for threshold tuning and provides robust regularization regardless of excitation level. We note this as a limitation of binary excitation gating in feedback-controlled marine vehicles, and recommend continuous regularization as a more robust alternative.
